\begin{figure}[t]
\centering
\begin{tikzpicture}
\begin{axis}[edge, height=0.6\linewidth,
    xmode=log, xlabel={Population, millions (log scale)},
    xmin=0.005, xmax=2000,
    xtick={0.01,0.1,1,10,100,1000}, xticklabels={0.01,0.1,1,10,100,1000},
    ymode=log, ymin=0.8, ymax=80,
    ytick={1,2,3,6,11,21,51}, yticklabels={0,1,2,5,10,20,50},
    ylabel={International submarine cables in service},
    legend pos=north west]
  \addplot[dotsBase, only if={other}{1}] table[x=pop_m,y=cables_p1] {\datadir/plot_countries.csv};
  \addlegendentry{Other coastal economies}
  \addplot[dotsSIDS, only if={sids}{1}] table[x=pop_m,y=cables_p1] {\datadir/plot_countries.csv};
  \addlegendentry{Small island states}
  \addplot[labelled=lab_cables] table[x=pop_m,y=cables_p1] {\datadir/plot_countries.csv};
\end{axis}
\end{tikzpicture}
\caption{\takeaway{Redundancy scales with size, not distance.} Cable counts
  rise with population (correlation of logs \corrCablesPop) but barely vary with
  remoteness (\corrCablesRemote). \nSIDSOneCable{} of \nSIDScables{} small
  island states, including \nPacificOneCable{} in the Pacific, have one
  international cable or none: a single cut takes them offline.
  \source{TeleGeography Submarine Cable Map (cables in service by 2025, excluding
  domestic systems); World Bank WDI.}}
\label{fig:cables}
\end{figure}
